\setchaptertoc
\chapter{Background}\label{chp:background}

\begin{summary}
    Use this summary environment to give short summaries at the beginning of each chapter.
    In~\Cref{sec:background:A}, we \dots.
    In ~\Cref{sec:background:A}, we \dots
\end{summary}

Use this chapter to introduce relevant parts of your topic in a way so that the average computer science reader understands it.

\section{Background Area A}\label{sec:background:A}
Descriptions, definitions, examples, and figures

\section{Background Area B}\label{sec:background:B}

\aiquote{GPT-5.2 other}{
    Generate a text snippets that symbolizes AI usage based on the ISF AI guidelines.
}{
    Worked quite well!
}{
    This paragraph symbolizes text generated with a type 2 AI tool that you want to use in your thesis. Please refer to the file \texttt{segments/ai-entries.tex} for section-wide and type 1 AI disclaimers.
}